\begin{abstract}
EEG foundation models aim to decode brain states across clinical and cognitive tasks from few
labels. They are pretrained on large unlabeled corpora with objectives borrowed from images and
language, reconstructing the raw waveform of masked patches or predicting discrete tokens, and
their progress has come mainly from more data and larger models. These objectives also fit
waveform details beyond the spectral information relevant to many EEG labels. Sleep stages,
clinical slowing and many cognitive states are associated with the power of neural rhythms,
while a short patch also contains phase and sample-level variation that may be only partly
predictable from context. We present \textbf{SpecBraM}, an EEG foundation model pretrained by
\emph{masked band-power prediction}: from the visible context, it predicts the log power of each
masked (channel, time) patch in the classical frequency bands, a parameter-free physical target
computed from the patch's spectrum. Under an idealized spectral model, band-averaged power is
sufficient for the band levels, and its logarithm stabilizes sampling variance when the numerical
floor is ignored. This motivates a compact target for efficient learning. With backbone, data
and training steps fixed, the trained models explain substantially more of the band-power
target's variance than of the waveform's, and on sleep staging band-power pretraining reaches
the linear-probe quality of waveform reconstruction with a fraction of the steps and with a
model a tenth of the size. Under an identical strict linear probe on eight datasets spanning six task types,
SpecBraM, pretrained on a fraction of the data of recent models, outperforms CBraMod, LaBraM
and CSBrain on nearly all tasks and has higher point estimates than REVE, which uses far more
data, on most tasks. On sleep staging it also leads the compared models with few labelled
training epochs and in unsupervised clustering agreement.
\end{abstract}
